\chapter{Sistemas monetarios y financieros internacionales}

\section{Introducción}

Los cuatro temas anteriores tratan el comercio de bienes y la política que lo
regula dando por supuesto que existe una forma de pagar lo que se compra fuera.
Este tema levanta ese supuesto. Cada Estado emite su propia moneda y no hay una
moneda internacional de curso legal, así que toda transacción con el exterior
—una importación, una inversión directa, la compra de un bono— exige antes
cambiar una moneda por otra. El precio al que se produce ese cambio, el mercado
donde se fija y las reglas que los países acuerdan para gobernarlo son el objeto
de este capítulo.

El tema se organiza en cuatro bloques. El primero define el tipo de cambio y
describe el mercado donde se determina. El segundo compara los dos regímenes
extremos, tipos fijos y tipos flexibles, y las fórmulas intermedias que se usan
en la práctica. El tercero aborda la integración financiera internacional y las
formas de medirla. El cuarto recorre la evolución del sistema monetario
internacional desde el patrón oro hasta el orden actual, que buena parte de la
literatura describe como la ausencia de un sistema.

El tratamiento de los tipos de cambio cruzados y de los ejercicios de arbitraje
sigue el planteamiento de \textcite{bengochea2009}; la parte histórica se apoya
en \textcite{requeijo2021} y la formalización de las condiciones de paridad, en
\textcite{feenstra2021macro}.

\section{El tipo de cambio y el mercado de divisas}

\subsection{Definición y formas de cotización}

El \textbf{tipo de cambio bilateral} entre dos monedas es el precio de una
expresado en unidades de la otra. Al ser un precio relativo, su valor depende de
cuál de las dos se tome como referencia, y de ahí que convivan dos convenciones
de cotización que dicen lo mismo del revés.

\begin{itemize}
    \item El \textbf{método directo}, llamado también americano, expresa el tipo
    de cambio como el número de unidades de moneda nacional necesarias para
    comprar una unidad de moneda extranjera. Bajo esta convención, un tipo
    euro/dólar de 1,12 significa que hacen falta 1,12~\euro{} para obtener un
    dólar.
    \item El \textbf{método indirecto}, o europeo, invierte la relación y expresa
    cuántas unidades de moneda extranjera se obtienen por una unidad de moneda
    nacional. El mismo caso se cotizaría como 0,89~\$ por euro.
\end{itemize}

El Banco Central Europeo publica sus tipos con el método indirecto, mientras que
la mayor parte de los manuales trabaja con el directo. La elección es
irrelevante siempre que se mantenga a lo largo de todo el análisis, porque una
convención es la inversa de la otra. Formalmente, si $TC_{ks}$ denota el tipo de
cambio de la moneda $k$ respecto de la $s$, se cumple la llamada
\textbf{condición de consistencia}:

\[
TC_{ks} = \frac{1}{TC_{sk}}
\qquad\Longrightarrow\qquad
TC_{ks} \times TC_{sk} = 1
\]

Con el método directo, un \textbf{aumento} del tipo de cambio significa que hay
que entregar más moneda nacional por la misma cantidad de moneda extranjera: la
moneda nacional se \textbf{deprecia}. Una \textbf{reducción} indica lo contrario
y la moneda se \textbf{aprecia}. Conviene reservar los términos
\emph{devaluación} y \emph{revaluación} para los movimientos equivalentes
decididos por la autoridad monetaria dentro de un régimen de tipos fijos, porque
la diferencia no es terminológica sino de mecanismo: en un caso el movimiento lo
produce el mercado y en el otro, una decisión de política económica.

\subsection{Tipos de cambio cruzados y arbitraje}

Cuando dos monedas no cotizan directamente entre sí, su tipo de cambio se deduce
a partir de las cotizaciones de ambas frente a una tercera. Ese tipo deducido es
el \textbf{tipo de cambio cruzado}:

\[
TC_{ks} = TC_{kj} \times TC_{js}
\]

La igualdad no es una definición contable sino una condición de equilibrio que
impone el mercado. Si el tipo cotizado se separa del cruzado, aparece una
oportunidad de \textbf{arbitraje triangular}: comprar la moneda infravalorada,
encadenar las tres operaciones y cerrar la posición con beneficio sin asumir
riesgo. Como esas operaciones se ejecutan en segundos y en volúmenes muy
grandes, la propia actividad arbitrajista devuelve las cotizaciones a la
consistencia. El interés del cálculo, en la práctica, no está en obtener el
beneficio sino en comprobar que las cotizaciones observadas son coherentes entre
sí.

\subsection{Del tipo nominal al tipo real y al tipo efectivo}

El tipo de cambio nominal compara monedas, no capacidad de compra. Para saber si
un país gana o pierde competitividad hay que corregirlo por los precios
relativos, lo que da el \textbf{tipo de cambio real}:

\[
TCR = TCN \times \frac{P^{*}}{P}
\]

donde $P^{*}$ es el nivel de precios extranjero y $P$ el nacional. Una
depreciación nominal que venga acompañada de una inflación interna equivalente
deja el tipo real inalterado y no mejora la posición competitiva: el efecto se
lo come el diferencial de precios.

Cuando se quiere resumir la posición de una moneda frente a un conjunto de
socios y no frente a uno solo, se construye el \textbf{tipo de cambio efectivo},
una media de los tipos bilaterales ponderada por el peso de cada socio en el
comercio del país:

\[
TCE_i = \left[ \sum_j \left( \frac{TC_{ij}}{TC_{i0}} \right) \times p_j \right]
\times 100
\]

con $TC_{i0}$ el tipo del período base y $p_j$ la ponderación del socio $j$. Si
además se corrige por precios se obtiene el tipo de cambio efectivo real, que es
el indicador que suelen seguir los organismos internacionales para valorar la
competitividad de una economía.

\subsection{La estructura del mercado de divisas}

El mercado de divisas no tiene sede física: es una red de operadores conectados
electrónicamente que funciona de forma continua. Sus rasgos relevantes son
cuatro.

\begin{enumerate}
    \item \textbf{Participantes.} Bancos comerciales, que concentran la mayor
    parte del volumen; empresas no financieras, que acuden por necesidades de
    cobro y pago derivadas del comercio; instituciones financieras no bancarias;
    y bancos centrales, cuya intervención es pequeña en volumen pero decisiva en
    señal.
    \item \textbf{Segmento al contado.} La entrega se produce en un plazo muy
    corto, habitualmente dos días hábiles, al tipo de cambio \emph{spot}.
    \item \textbf{Segmento a plazo.} Comprador y vendedor fijan hoy el tipo al
    que intercambiarán las monedas en una fecha futura. El tipo \emph{forward}
    resultante no es una previsión del tipo futuro, sino el que hace indiferente
    invertir en una u otra moneda dado el diferencial de tipos de interés.
    \item \textbf{Función económica.} Más allá de liquidar transacciones, el
    mercado permite \textbf{cubrir riesgo} de tipo de cambio y también
    \textbf{especular} con él, según se tome la posición para compensar una
    exposición existente o para crearla.
\end{enumerate}

\section{Los regímenes cambiarios}

\subsection{Tipos de cambio fijos}

En un régimen de tipos fijos la autoridad monetaria se compromete a mantener la
cotización de su moneda en un valor determinado, o dentro de una banda estrecha
alrededor de él, y a intervenir comprando o vendiendo divisas cuando el mercado
tienda a separarse de ese valor.

La ventaja principal es la \textbf{certidumbre}: elimina el riesgo de tipo de
cambio en las operaciones con los países del área, lo que favorece el comercio y
la inversión, y proporciona un ancla nominal que disciplina la inflación en
economías con historial de precios inestables. El coste es la
\textbf{pérdida de autonomía monetaria}. Defender la paridad obliga a supeditar
el tipo de interés a esa defensa, de modo que la política monetaria deja de
poder usarse para estabilizar la producción y el empleo. A ello se suma la
necesidad de mantener reservas suficientes y la exposición a ataques
especulativos cuando el mercado percibe que la paridad no es sostenible.

\subsection{Tipos de cambio flexibles}

En un régimen flexible el tipo lo determina la oferta y la demanda de divisas
sin compromiso de intervención. La ventaja simétrica de la anterior es la
\textbf{autonomía}: el banco central puede fijar el tipo de interés atendiendo a
la situación interna, y el tipo de cambio actúa como amortiguador automático
ante perturbaciones externas, ya que una caída de la demanda exterior deprecia
la moneda y ese abaratamiento compensa en parte el golpe.

El inconveniente es la \textbf{volatilidad}. Los tipos flexibles se mueven mucho
más de lo que justifican los fundamentos, en parte por el fenómeno de
\emph{overshooting}: como los precios de los bienes son rígidos a corto plazo y
los de los activos no, un cambio en la política monetaria hace que el tipo de
cambio se desplace inicialmente más allá de su valor de largo plazo y vuelva
después. Esa volatilidad encarece la cobertura y añade incertidumbre a las
decisiones de comercio e inversión.

\subsection{Los regímenes intermedios}

Entre los dos extremos hay un abanico de fórmulas que combinan rasgos de ambos:
bandas de fluctuación, paridades móviles o \emph{crawling peg}, flotación
sucia—en la que se deja flotar la moneda pero el banco central interviene de
forma discrecional—, cajas de conversión o \emph{currency board}, que respaldan
cada unidad emitida con reservas y renuncian de hecho a la política monetaria, y
la dolarización o adopción directa de una moneda ajena.

\subsection{El trilema de la economía abierta}

La elección de régimen no es libre: está sujeta a una restricción conocida como
\textbf{trilema} o triángulo de imposibilidad. De estos tres objetivos, solo se
pueden alcanzar dos a la vez:

\begin{enumerate}
    \item tipo de cambio fijo,
    \item libertad de movimientos de capital,
    \item política monetaria autónoma.
\end{enumerate}

El patrón oro renunciaba al tercero. Bretton Woods conservó los dos primeros
limitando la movilidad de capital. La eurozona resuelve el trilema cediendo la
política monetaria a una autoridad común. Y una economía con tipo flexible y
capital libre, como la mayoría de las desarrolladas hoy, renuncia al primero. El
trilema explica buena parte de la historia que sigue: cada vez que un país ha
intentado los tres a la vez, el mercado le ha obligado a elegir.

\section{La integración financiera internacional}

La integración financiera internacional mide hasta qué punto los mercados de
capitales de distintos países funcionan como uno solo. Su relevancia es doble:
permite que el ahorro se dirija a donde es más productivo y que los países
suavicen su consumo ante perturbaciones transitorias, pero también transmite las
crisis con rapidez y expone a las economías a reversiones bruscas de los flujos.

\subsection{Las condiciones de paridad}

La forma más habitual de medir la integración es comprobar si se cumplen las
condiciones de arbitraje que la competencia perfecta impondría.

La \textbf{paridad cubierta de tipos de interés} compara invertir en el país con
hacerlo fuera cubriendo el riesgo de cambio en el mercado a plazo:

\[
1 + r = (1 + r^{*}) \times \frac{F}{E}
\]

donde $r$ y $r^{*}$ son los tipos de interés nacional y extranjero, $E$ el tipo
al contado y $F$ el tipo a plazo. Su incumplimiento sistemático señala barreras
efectivas a los movimientos de capital, porque la operación no tiene riesgo.

La \textbf{paridad descubierta} sustituye el tipo a plazo por la expectativa de
tipo futuro y, en consecuencia, incorpora riesgo:

\[
r = r^{*} + \frac{\Delta E^{e}}{E} + \rho
\]

El término $\rho$ recoge la prima de riesgo. Su utilidad práctica se ve bien en
episodios de tensión: cuando la prima de riesgo de la deuda de un país aumenta,
los inversores exigen más rentabilidad, venden sus bonos y el precio cae, con lo
que la propia prima se realimenta.

\subsection{Medidas de cantidad}

Junto a las de precios se emplean medidas basadas en volúmenes: el peso de los
activos y pasivos exteriores sobre el PIB, el tamaño de los flujos brutos de
capital y la correlación entre ahorro e inversión nacionales. Esta última es la
base de la \textbf{paradoja de Feldstein-Horioka}: si el capital fuese
perfectamente móvil, la inversión de un país no tendría por qué financiarse con
su propio ahorro y ambas magnitudes deberían estar poco correlacionadas. La
evidencia muestra una correlación alta incluso entre economías abiertas, lo que
sugiere que la integración es menor de lo que los datos de flujos sugieren.

\section{La evolución del sistema monetario internacional}

Se entiende por \textbf{sistema monetario internacional} el conjunto de reglas,
instituciones y prácticas que gobiernan los pagos entre países: cómo se
determinan los tipos de cambio, qué activos sirven como reserva y cómo se
corrigen los desequilibrios de balanza de pagos. La historia de los últimos
ciento cincuenta años puede leerse como una sucesión de respuestas distintas a
esas tres preguntas.

\subsection{El patrón oro}

Vigente aproximadamente entre 1870 y 1914, el patrón oro fijaba el valor de cada
moneda en una cantidad determinada de metal y garantizaba su convertibilidad.
Los tipos de cambio quedaban así fijados de hecho, y el ajuste de los
desequilibrios era automático a través del mecanismo de flujo de oro y precios:
un país con déficit perdía oro, su masa monetaria y sus precios caían, ganaba
competitividad y el déficit se corregía.

El sistema funcionó con notable estabilidad, pero conviene no atribuir todo el
mérito al mecanismo. Coincidió con un período de paz y de ciclos económicos
sincronizados, y contó con el liderazgo británico y con un grado apreciable de
cooperación entre los bancos centrales del núcleo, que se prestaban oro cuando
alguno atravesaba dificultades. El coste era el que impone el trilema: sin
política monetaria propia, el ajuste recaía íntegramente sobre los precios y el
empleo.

\subsection{El período de entreguerras}

La Primera Guerra Mundial interrumpió la convertibilidad. Los intentos de volver
al patrón oro en los años veinte fracasaron, en buena medida porque algunos
países regresaron con paridades que ya no correspondían a sus precios: el caso
británico de 1925, con una libra sobrevalorada, es el ejemplo habitual. La
Depresión terminó de romper el orden y dio paso a devaluaciones competitivas,
control de cambios y proteccionismo. Es el período que las conferencias de
posguerra tuvieron presente al diseñar lo que vino después.

\subsection{El sistema de Bretton Woods}

El acuerdo de 1944 estableció un sistema de \textbf{tipos de cambio fijos pero
ajustables}: cada moneda se definía frente al dólar, y el dólar era convertible
en oro a 35 dólares la onza. Los países se comprometían a mantener su paridad
dentro de una banda estrecha y podían modificarla solo ante un desequilibrio
fundamental, con autorización del recién creado Fondo Monetario Internacional.
El FMI proporcionaba además financiación para desequilibrios transitorios, y la
movilidad de capital quedaba limitada, que es lo que permitía conservar a la vez
tipos fijos y cierta autonomía monetaria.

\subsection{La crisis y el final de Bretton Woods}

El sistema llevaba dentro una contradicción que Robert Triffin formuló a finales
de los cincuenta y que se conoce como \textbf{dilema de Triffin}: la liquidez
internacional dependía de que Estados Unidos suministrase dólares mediante
déficits, pero cuanto mayor era el volumen de dólares en circulación frente a
las reservas de oro, menos creíble resultaba el compromiso de convertibilidad.

La creación de los \textbf{Derechos Especiales de Giro} en 1969 fue un intento
tardío de dotar al sistema de un activo de reserva que no dependiera de una
moneda nacional. No bastó. En agosto de 1971 Estados Unidos suspendió la
convertibilidad del dólar en oro y, tras un intento de reajuste de paridades en
los Acuerdos Smithsonianos, el sistema se abandonó definitivamente en 1973.

\subsection{Los euromercados}

En paralelo se había desarrollado un mercado de depósitos y préstamos
denominados en una moneda distinta de la del país donde se realiza la operación:
los eurodólares y, en general, los euromercados. Su origen está en la búsqueda
de intermediación al margen de la regulación y del control de cambios
nacionales. Su crecimiento importa aquí por dos razones: aumentó de forma
notable la movilidad efectiva del capital, erosionando el supuesto sobre el que
se sostenía Bretton Woods, y creó un canal de financiación internacional que
después resultaría decisivo en el reciclaje de los excedentes petrolíferos de
los años setenta.

\subsection{El sistema monetario internacional después de 1973}

Los acuerdos de Jamaica de 1976 reconocieron formalmente lo que ya ocurría: cada
país es libre de elegir su régimen cambiario. El resultado es un orden híbrido
en el que conviven monedas en flotación, monedas ancladas a otra y uniones
monetarias, con el dólar conservando el papel principal como moneda de reserva y
de facturación pese a no tener ya respaldo metálico. El FMI cambió de función:
de garante de paridades pasó a prestamista y supervisor, papel que ejerció en
las sucesivas crisis de deuda y de balanza de pagos de las décadas siguientes.

\subsection{Del Sistema Monetario Europeo a la Unión Económica y Monetaria}

Europa respondió a la inestabilidad posterior a 1973 con acuerdos propios de
estabilidad cambiaria. La \emph{serpiente monetaria} dio paso en 1979 al
\textbf{Sistema Monetario Europeo}, articulado sobre el ECU como unidad de
cuenta y un mecanismo de tipos de cambio con bandas de fluctuación y obligación
de intervención. La crisis de 1992 y 1993 obligó a ampliar las bandas y mostró
la dificultad de sostener paridades con capital libre.

El Tratado de Maastricht fijó el camino hacia la \textbf{Unión Económica y
Monetaria} y los criterios de convergencia en inflación, tipos de interés,
déficit, deuda y estabilidad cambiaria. El euro se introdujo como moneda
escritural en 1999 y en circulación en 2002, con el Banco Central Europeo como
autoridad monetaria única. La discusión académica sobre si la eurozona
constituye un \textbf{área monetaria óptima} en el sentido de Mundell —movilidad
del trabajo, flexibilidad de precios y salarios, integración fiscal y
sincronía de los ciclos— sigue abierta, y la crisis de deuda soberana la
reactivó al mostrar que la unión monetaria se había construido sin una unión
fiscal que la acompañase \parencite{cuenca2019}.

\subsection{El «no sistema» y el debate sobre su reforma}

Buena parte de la literatura describe el orden actual como un \textbf{no
sistema}: no hay reglas comunes sobre regímenes cambiarios, ni un mecanismo
acordado de ajuste de los desequilibrios, ni un activo de reserva supranacional
con peso real. Las propuestas de reforma discutidas desde los años setenta van
en tres direcciones: recuperar alguna forma de zonas objetivo para los tipos de
cambio, ampliar el papel de los Derechos Especiales de Giro como activo de
reserva, y gravar las transacciones financieras de corto plazo para reducir la
volatilidad, en la línea de la propuesta de Tobin. Ninguna se ha adoptado, y el
debate sobre la arquitectura financiera internacional continúa
\parencite{requeijo2021,alonso2021}.
